% !TeX root = ../main.tex
\section{Phân tích tính khả thi và mô hình phí}

\subsection{Tính khả thi về thị trường}

Mutux có tính khả thi vì thị trường game và eSports tại Việt Nam đang có nhiều tín hiệu phát triển tích cực. Khi game được nhìn nhận như một ngành công nghiệp văn hóa trọng điểm, các sản phẩm phục vụ hệ sinh thái gaming cũng có thêm cơ hội phát triển, trong đó có gaming gear như chuột, bàn phím cơ, tai nghe, tay cầm và phụ kiện eSports.

Nhu cầu này không chỉ đến từ việc chơi game giải trí, mà còn đến từ nhóm người chơi nghiêm túc hơn, muốn chọn thiết bị phù hợp với kỹ năng, thói quen chơi và ngân sách. Với nhóm này, thông số kỹ thuật hoặc review trên mạng chưa đủ để ra quyết định mua, vì cảm giác sử dụng thực tế mới là yếu tố quan trọng.

\subsection{Tính khả thi về nhu cầu người dùng}

Người dùng có lý do rõ ràng để sử dụng Mutux: họ muốn thử gear thật trước khi mua, nhưng không muốn bỏ ra một khoản tiền lớn để mua nhầm thiết bị. Một sản phẩm gaming gear tốt có thể có giá từ vài trăm nghìn đến vài triệu đồng. Nếu mua sai, người dùng vừa mất tiền, vừa khó bán lại với giá tốt.

Mức phí thuê ngắn ngày có thể dễ chấp nhận hơn nhiều so với chi phí mua mới. Ví dụ, người dùng có thể trả khoảng 30.000--100.000 đồng mỗi ngày, hoặc khoảng 150.000 đồng cho ba ngày thuê một thiết bị. Khoản phí này hợp lý hơn so với việc mua ngay một thiết bị trị giá 1--3 triệu đồng khi chưa chắc nó có phù hợp hay không.

Tuy nhiên, nếu chỉ có phí thuê thì mô hình vẫn chưa đủ hấp dẫn, vì người thuê thường phải đặt cọc cao để bảo vệ tài sản cho shop. Trong cách thuê truyền thống, khoản cọc có thể bằng toàn bộ giá trị sản phẩm, nên một món hàng dưới 2.000.000 đồng vẫn khiến người thuê phải chuẩn bị gần đủ số tiền đó cộng thêm tiền thuê. Đây là rào cản lớn, đặc biệt với sinh viên, người mới chơi eSports hoặc người có ngân sách hạn chế.

\subsection{Cách hoạt động và các khoản phí}

Mutux hoạt động theo mô hình marketplace kết nối ba bên chính: người thuê gear, shop hoặc chủ gear, và nền tảng Mutux. Trong phiên bản mở rộng, Mutux có thể kết nối thêm đối tác tài chính để triển khai dịch vụ Cọc Trả Sau, tức hỗ trợ vay hoặc bảo chứng phần cọc theo giá trị sản phẩm.

Các khoản phí chính trong mô hình gồm:

\begin{itemize}
    \item \textbf{Phí thuê gear}: người thuê trả theo ngày hoặc theo gói ngắn hạn. Ví dụ, một tay cầm trị giá khoảng 2.000.000 đồng có thể được cho thuê với giá 10.000 đồng/ngày; nếu thuê ba ngày thì phí thuê là 30.000 đồng.
    \item \textbf{Hoa hồng nền tảng}: Mutux lấy khoảng 10\%--20\% trên phí thuê thành công. Nếu phí thuê là 30.000 đồng, nền tảng có thể nhận 3.000--6.000 đồng, phần còn lại thuộc về shop hoặc chủ gear.
    \item \textbf{Tiền cọc truyền thống}: người thuê đặt cọc để bảo vệ tài sản cho bên cho thuê. Với thiết bị trị giá khoảng 1.800.000--2.000.000 đồng, tiền cọc có thể bằng toàn bộ giá trị sản phẩm. Khoản này không phải doanh thu của Mutux mà là khoản đảm bảo, được hoàn lại khi người thuê trả thiết bị đúng tình trạng.
    \item \textbf{Phí Cọc Trả Sau hoặc phí vay cọc}: nếu người dùng không muốn hoặc không đủ khả năng ứng trước toàn bộ tiền cọc, họ có thể chọn dịch vụ hỗ trợ cọc. Ví dụ, thay vì khóa trước 2.000.000 đồng tiền cọc cho một món gear, người dùng trả thêm khoảng 5.000--10.000 đồng phí vay/bảo chứng cọc cho một đơn thuê ngắn ngày. Mức phí này được thiết kế rất thấp, tương tự các nền tảng ví trả sau hoặc SPayLater, còn tiền thuê/mượn thiết bị vẫn do người cho thuê đặt ra.
    \item \textbf{Hoa hồng từ đối tác tài chính}: nếu dịch vụ Cọc Trả Sau được triển khai qua ngân hàng, ví điện tử hoặc công ty tài chính được cấp phép, Mutux có thể nhận phần chia sẻ doanh thu hoặc hoa hồng giới thiệu từ giao dịch này.
    \item \textbf{Hoa hồng mua sau thuê}: nếu người dùng thuê thử và quyết định mua sản phẩm từ shop, Mutux có thể nhận hoa hồng từ giao dịch mua. Shop cũng có thể trừ một phần phí thuê vào giá bán để tăng tỷ lệ chuyển đổi.
    \item \textbf{Gói thành viên}: người dùng trả phí định kỳ để được giảm phí thuê, giảm phí bảo lãnh cọc, ưu tiên thuê gear hot hoặc nhận ưu đãi giao nhận.
    \item \textbf{Gói dịch vụ cho shop}: shop có thể trả phí để đăng nhiều thiết bị hơn, dùng dashboard quản lý đơn thuê, được quảng bá sản phẩm hoặc tiếp cận nhóm khách có nhu cầu mua sau khi dùng thử.
\end{itemize}

Với cách vận hành này, doanh thu của Mutux không phụ thuộc hoàn toàn vào phí thuê gear. Phần phí thuê chủ yếu cần đủ hấp dẫn để shop và chủ gear có động lực đưa thiết bị lên nền tảng. Lớp doanh thu quan trọng hơn nằm ở hoa hồng, phí vay/hỗ trợ cọc, chia sẻ doanh thu với đối tác tài chính và giao dịch mua sau thuê.

\subsection{So sánh hai cách đặt cọc}

Điểm khác biệt quan trọng của Mutux là cho người dùng lựa chọn giữa đặt cọc truyền thống và Cọc Trả Sau. Ví dụ, người dùng thuê một tay cầm trị giá 2.000.000 đồng trong ba ngày, với phí thuê 10.000 đồng/ngày.

\begin{center}
\begin{tabular}{>{\raggedright\arraybackslash}p{0.36\textwidth} >{\raggedleft\arraybackslash}p{0.22\textwidth} >{\raggedleft\arraybackslash}p{0.24\textwidth}}
\toprule
\textbf{Khoản mục} & \textbf{Cọc truyền thống} & \textbf{Cọc Trả Sau} \\
\midrule
Giá trị tay cầm & 2.000.000đ & 2.000.000đ \\
Phí thuê 3 ngày & 30.000đ & 30.000đ \\
Tiền cọc ứng trước & 2.000.000đ & 0đ \\
Phí vay/hỗ trợ cọc & 0đ & 5.000--10.000đ \\
Tổng tiền ban đầu & 2.030.000đ & 35.000--40.000đ \\
\bottomrule
\end{tabular}
\end{center}

Như vậy, với cùng một món tay cầm trị giá 2.000.000 đồng, cách thuê truyền thống khiến người dùng phải chuẩn bị 2.030.000 đồng ngay từ đầu, còn Cọc Trả Sau chỉ cần khoảng 35.000--40.000 đồng cho ba ngày thuê. Khoản chênh lệch này làm mô hình đủ tiện và hợp lý hơn cho sinh viên, người mới chơi eSports hoặc người có ngân sách hạn chế: họ vẫn được dùng thử thiết bị thật, bên cho thuê vẫn có cơ chế bảo vệ tài sản nếu hư hỏng hoặc mất đồ, còn Mutux có thêm nguồn thu ngoài hoa hồng thuê.

\subsection{Tính khả thi về vận hành}

Ở giai đoạn MVP, Mutux chưa cần tự đứng ra cho vay hoặc triển khai sản phẩm tài chính thật. Hướng an toàn hơn là mô phỏng flow Cọc Trả Sau để chứng minh nhu cầu, sau đó mới hợp tác với ngân hàng, ví điện tử hoặc công ty tài chính được cấp phép nếu muốn triển khai thật.

Để mô hình vận hành được, nền tảng cần có các cơ chế cơ bản:

\begin{itemize}
    \item Xác minh danh tính người thuê trước khi cho thuê thiết bị có giá trị cao.
    \item Kiểm tra và ghi nhận tình trạng thiết bị trước và sau khi thuê.
    \item Quy định rõ phí trả muộn, phí hư hỏng, mất thiết bị và điều kiện hoàn cọc.
    \item Theo dõi trạng thái thiết bị: available, booked, renting, returned, maintenance.
    \item Giới hạn giá trị hỗ trợ cọc theo độ tin cậy của người thuê.
    \item Ưu tiên thử nghiệm với một số shop hoặc cộng đồng gamer nhỏ trước khi mở rộng.
\end{itemize}

Nhờ đó, Mutux có thể kiểm chứng nhu cầu thị trường trước, đồng thời tránh rủi ro pháp lý khi chưa đủ điều kiện tự cung cấp dịch vụ tài chính.
